\section{The complete quadratic seven-factor family on cubic coefficients}
\label{sec:mixed-cubic-quadratic-seven-strict}

Put $K=\mathbb Q(\omega)$, $t^3=269$, $L=K(t)$, and let $\chi$
be either faithful character of $L/K$.  Let $\xi_7$ be a faithful
rational character of conductor seven and order six, valued in the
original coefficient ring.  Use the exact dictionary
\[
 \begin{array}{c|rrrr}
  d&-7&7&-14&14\\\hline
  \beta_d&1&\chi_{-4}&\chi_8&\chi_{-8}\\
  \delta_d=\beta_d\xi_7^3&\chi_{-7}&\chi_{28}&\chi_{-56}&\chi_{56}.
 \end{array}
 \tag{MQS.1}
\]
Thus $\delta_d$ is the primitive quadratic character of
$\mathbb Q(\sqrt d)$; every character is restricted to $G_K$.
For $d$ in this four-element set put
\[
 B_d=\mathbb Q(t,\sqrt d),\qquad F_d=L(\sqrt d),\qquad
          s=\sqrt d,\quad u=t+s.
 \tag{MQS.2}
\]
Let $g$ send $t$ to $\omega t$ and fix $s$, let $b$ fix
$t,\omega$ and negate $s$, and let $\tau$ fix $t,s$ and
invert $\omega$.

The completed finite computations proved in Step1 give
\[
 \begin{gathered}
  (m_{-7},m_7,m_{-14},m_{14})=(21,7,28,28),\\
  C_{m_d}\twoheadrightarrow\operatorname{Cl}(B_d),\qquad
       b=1\quad\hbox{on }\operatorname{Cl}(B_d)\{7\}.
 \end{gathered}
 \tag{MQS.3}
\]
The first map is the natural class map of the displayed ideal
generator.  The second assertion follows from exact principal
identities for its seven-primary power.
No equality of actual class numbers or complete unit basis is needed.

For an odd coefficient prime $\ell>3$, let $O$ be the original
ring of integers of any finite $E/\mathbb Q_\ell$ containing
$\mu_3$, with original uniformizer $\varpi$, residue field $k$
and embedding.  Let $\kappa:G_K\to1+\mathfrak m_O$ be continuous
and unramified away $\ell$.  For $j=1,2$, $0\leq a\leq5$
and a positive integer $e$, put
\[
         D=(E/O)(\chi^j\beta_d\xi_7^a\kappa^{-1}c_\ell^{-e}).
 \tag{MQS.4}
\]
Full strict means zero restriction to $H^1(K_v,D)$ at every
$v\mid\ell$, and to $H^1(I_v,D)$ at every other finite place.
Use exactly these conditions for $D[\varpi]$.

Define for each $d$ the split-prime set
\[
 \begin{split}
 \mathscr P_d=\{\ell\text{ prime}:{}&\ell>3,\quad\ell\ne7,269,
       \quad\ell\equiv1\pmod3,\\
       &269^{(\ell-1)/3}\not\equiv1\pmod\ell
                            \ \text{or}\ \delta_d(\ell)=-1\}.
 \end{split}
 \tag{MQS.5}
\]

\begin{theorem}[AC-G217: the full quadratic seven-factor family]
\label{AC-G217}
The following entire family has full original-coefficient strict zero:
\[
                  \mathcal S_\ell(K,D[\varpi])=0,
                  \qquad\mathcal S_\ell(K,D)=0.
 \tag{MQS.6}
\]
First, take any $d$, $\ell\in\mathscr P_d$, $j=1,2$, $a=3$,
and positive $e$ divisible by $\ell-1$.
Second, take coefficient seven, positive $e$, $j=1,2$, and
any of the original six $a$ for which the literal rational
residual character
\[
                    \rho_{a,e}=\bar\xi_7^{\,a}\bar c_7^{-e}
 \tag{MQS.7}
\]
satisfies either $\operatorname{ord}(\rho_{a,e})=2$ with any $d$,
or $\rho_{a,e}=1$ with $\beta_d\ne1$.
The full finite character $\xi_7^a$ in this part may have
order six; only its product with the reduced cyclotomic twist is
quadratic or trivial.  The residual-trivial case uses the exact
arbitrary-near-one-character scope of AC-G211, without a new
class computation or a coefficient change.

For every integer $T\geq2$, put $N=T-1$ and define
\[
 \begin{split}
 \mathcal V_u(T)=\{(\ell,\chi^j\beta_d\xi_7^3,e):{}&d\in\{-7,7,-14,14\},
       \ \ell\in\mathscr P_d,\ j=1,2,\\
                         &1\leq e<T,\quad\ell-1\mid e\},\\
 \mathcal V_7(T)=\{(7,\chi^j\beta_d\xi_7^a,e):{}&d\in\{-7,7,-14,14\},
       \ j=1,2,\ 0\leq a\leq5,\\
                         &1\leq e<T,\quad
          [\operatorname{ord}(\rho_{a,e})=2
             \ \text{or}\ (\rho_{a,e}=1,\ \beta_d\ne1)]\}.
 \end{split}
 \tag{MQS.8}
\]
Their exact cardinalities are
\[
       |\mathcal V_u(T)|=
          2\sum_d\sum_{\substack{\ell\in\mathscr P_d\\\ell\leq T}}
                           \left\lfloor\frac{N}{\ell-1}\right\rfloor,
       \qquad |\mathcal V_7(T)|=14N.
 \tag{MQS.9}
\]
Their coefficient primes make the two sets disjoint.
The exact intersection of $\mathcal V_7(T)$ with the previously
removed families is the $6\lfloor N/6\rfloor$ AC-G211 triples
with $\beta_d\ne1$, $a=0$ and $6\mid e$.
Consequently the newly removed coefficient-seven set has
$14N-6\lfloor N/6\rfloor$ triples.  The unramified-prime set
is entirely new.  For every retained set $J$ of actual triples
in (CP.3),
\[
          \mathcal R_J=
           \mathcal R_{J\setminus(\mathcal V_u(T)\cup\mathcal V_7(T))}.
 \tag{MQS.10}
\]
All original coefficient orders, embeddings, ranks $b_\ell$,
completion and inertia kernels are unchanged.  This deletes full
mixed strict summands with zero cost and on the identical window;
it makes no individual positive-derived-index assertion.
\end{theorem}

\begin{proof}
\proofstep{1. Exact fields, automorphisms and natural class inputs.}
The splitting field of $X^3-269$ is $L/\mathbb Q$ with group
$S_3$ and unique quadratic subfield $K$.
None of the four $\mathbb Q(\sqrt d)$ equals $K$; hence its
intersection with $L$ is $\mathbb Q$.  Therefore
\[
 \begin{gathered}
   [B_d:\mathbb Q]=6,\quad [F_d:\mathbb Q]=12,\quad
   F_d^{\langle\tau\rangle}=B_d,\\
   \operatorname{Gal}(F_d/K)=\langle g\rangle\times\langle b\rangle
               \simeq C_3\times C_2,\qquad
   \tau g\tau^{-1}=g^{-1},\quad\tau b=b\tau.
 \end{gathered}
 \tag{MQS.11}
\]
The six sums $t_i\pm\sqrt d$ are distinct.  A collision
between unequal cubic roots puts their nonzero difference in
$L\cap\mathbb Q(\sqrt d)=\mathbb Q$; an $S_3$ transposition
negates that difference, a contradiction.  Thus $u=t+s$ is
primitive.  Its exact degree-six polynomial and inverse are
\[
 \begin{gathered}
  Q_d(X)=(X^3+3dX-269)^2-d(3X^2+d)^2\\
  \hphantom{Q_d(X)}=X^6-3dX^4-538X^3+3d^2X^2-1614dX+72361-d^3,\\
       s=\frac{u^3+3du-269}{3u^2+d},\qquad t=u-s,
                       \qquad b(u)=u-2s.
 \end{gathered}
 \tag{MQS.12}
\]
The denominator is nonzero by irreducibility and its smaller
degree.  Exact polynomial reduction checks $s^2=d$, $t^3=269$
and the stated automorphism.  These identify the class and ideal
calculations with the required field and quadratic action.

For each exact polynomial (MQS.12), the retained computation
performs \texttt{bnfinit} with flag one and then the distinct
\texttt{bnfcertify} flag-one quotient check.  The complete outputs are
\[
 \begin{array}{r|r|c|r|c}
  d&m_d&\text{signature}&\text{discriminant}&\text{certificate}\\\hline
  -7&21&[0,3]&-16163884908927&1\\
   7&7&[2,2]&1034488634171328&1\\
  -14&28&[0,3]&-8275909073370624&1\\
   14&28&[2,2]&8275909073370624&1.
 \end{array}
 \tag{MQS.12a}
\]
The $d=-7$ run finished at 10:31:19.324229 UTC on
23 November 2026; the other three finished at
10:34:55.103803 UTC.  Both processes exited with code zero,
all expected certificate markers and complete outputs.
The documented scope is the natural cyclic quotient in
(MQS.3), not class-number equality or a complete unit basis.

Let $\mathfrak g_d$ be the displayed Smith-generator ideal.
Since every $m_d$ has exactly one factor seven,
$\mathfrak I_d=\mathfrak g_d^{m_d/7}$ generates the actual
seven-primary class group, possibly the trivial group.
The exact action checks use powers $3,1,4,4$, respectively,
and prove $b(\mathfrak I_d)=(\alpha_d)\mathfrak I_d$.
All four algebraic numbers, integral bases and ideal lattices
are retained with every rational denominator.
For example the $d=-7$ calculation gives
\[
 \begin{aligned}
      [\alpha_{-7}]_{\mathrm{zk}}=
        \bigl(&-27069,-15017,-2837/2,\\
              &17067/4,-12003/4,8301/8\bigr),
              \qquad \alpha_7=1.
 \end{aligned}
 \tag{MQS.13}
\]
The first action process finished at 10:32:10.992654 UTC;
the remaining three finished at 10:35:53.396344 UTC.
They verify all twenty-four field/automorphism identities
and all four literal principal ideal-HNF equalities, with
exit code zero and no errors.  The principal solver only
proposes the algebraic numbers; the exact ideal equalities
verify them, rather than trusting a class-coordinate vector.
The natural class-quotient map sends the computed Smith generators
to these ideal classes, as specified by the applied PARI
factorbase/Smith description and flag-one prime test.
This proves the quadratic-action assertion in (MQS.3).
The independent fresh class and alternate-generator executions
are retained separately from these root-owned executions;
their rational-lattice verifiers have their own precise scope.

\proofstep{2. Only the required class sectors vanish.}
For $\ell\in\mathscr P_d$, the quotients in (MQS.3) give
$\operatorname{Cl}(B_d)\{\ell\}=0$.
Set $C=\operatorname{Cl}(F_d)\{\ell\}$ and use additive notation.
The exact norm/extension identities are
\[
      N j_B=2\quad\hbox{on }\operatorname{Cl}(B_d)\{\ell\},
      \qquad j_BN=1+\tau\quad\hbox{on }C.
 \tag{MQS.14}
\]
Thus $\tau=-1$ on $C$.  Conjugating $g$ by this scalar in
(MQS.11) gives $g=g^{-1}$; since $g^3=1$, it follows that
$g=1$ on $C$.  Every equivariant map $C\to k$ whose
$g$-character is the nontrivial $\bar\chi^j$ is zero.

At coefficient seven use $C=\operatorname{Cl}(F_d)\{7\}$ and
the integral projectors
\[
             C=C^+\oplus C^-,\qquad
             C^+=\frac{1+b}{2}C,\quad C^-=\frac{1-b}{2}C.
 \tag{MQS.15}
\]
By the quadratic action in (MQS.3), the norm to $B_d$ is zero on $C^-$.
Equation (MQS.14) therefore gives $\tau=-1$ on $C^-$.
Both $g$ and $\tau$ preserve this summand, and the same
conjugation argument gives $g=1$ on it.
An equivariant homomorphism $C\to k$ with nontrivial
$g$-character and quadratic $b$-character is zero on
$C^+$ by the latter and on $C^-$ by the former.
No entire primary class group of $F_d$ has been asserted zero.

\proofstep{3. The residual coefficient and full strict restriction.}
In the first part, $a=3$, $\ell-1\mid e$ and
$\bar\kappa=1$ give
$A=D[\varpi]=k(\bar\chi^j\bar\delta_d)$.
For the order-two case at coefficient seven, the unique nontrivial
quadratic rational conductor-seven character is $\bar\xi_7^3$.
Consequently (MQS.7) gives the identical residual module
\[
                         A=k(\bar\chi^j\bar\delta_d).
 \tag{MQS.16}
\]
It is trivial on $G_{F_d}$; the full coefficient need not
be trivial there.  The degree $[F_d:K]=6$ is prime to
every coefficient prime under consideration.
Inflation--restriction and averaging therefore inject
$H^1(K,A)$ into $H^1(F_d,k)=\operatorname{Hom}_{\rm cont}(G_{F_d},k)$.
A residual full strict class restricts to a character
unramified at every finite place: its completion
restrictions are zero above the coefficient prime,
and its inertia restrictions are zero elsewhere.
The Hilbert class-field identification gives a map
$f:\operatorname{Cl}(F_d)\to k$ with exact natural
conjugation equivariance
\[
             f(\sigma x\sigma^{-1})=
                 (\bar\chi^j\bar\delta_d)(\sigma) f(x).
 \tag{MQS.17}
\]
It is zero by the targeted class argument in Step2.
Thus $\mathcal S_\ell(K,A)=0$ in these two cases.

\proofstep{4. Every original coefficient-place invariant vanishes.}
In the first part $\ell$ splits in $K$.  If $269$ is a
noncube modulo $\ell$, the two Kummer Frobenius values
are inverse primitive cube roots.  Multiplication by
$\delta_d(\ell)\in\{1,-1\}$ cannot make either value one.
If $269$ is a cube, Hensel's lemma and
$\mu_3\subset\mathbb Q_\ell$ put all three roots locally
in $\mathbb Q_\ell$, so the Kummer character is trivial
at both places.  The admitted alternative
$\delta_d(\ell)=-1$ again gives a nontrivial residual
value at each.  In both cases a Frobenius lift has full
scalar action differing from one by an $O$-unit,
independently of the possibly ramified near-one $\kappa$.
Therefore $D^{G_{K_v}}=0$ at both coefficient places.

For the order-two case at coefficient seven the two completions
are $\mathbb Q_7$.
Each $\mathbb Q(\sqrt d)$ is ramified quadratic at seven;
$L/K$ and every $\beta_d$ are unramified there.
The admitted $\rho_{a,e}$ has order two and nontrivial
inertia on each completion.  Its value $-1$ differs
from one by a unit of $O$.  Thus the full scalar action
on some inertia element has unit difference at each
place, and again $D^{G_{K_v}}=0$.  This argument uses
the actual product $\bar\xi_7^a\bar c_7^{-e}$ and retains
all six possible original finite characters.

For clarity, no inert unramified-prime completion is
silently included in (MQS.5).  If $\ell\equiv2\pmod3$,
cubing is bijective on $\mathbb F_\ell^\times$; the root
lifts to $\mathbb Q_\ell$, and all roots lie in $K_v$.
The Kummer factor is locally trivial.  The rational
quadratic Frobenius is $\delta_d(\ell)^2=1$ because
$[k(v):\mathbb F_\ell]=2$.  The required local invariant
quotient is not removed by this argument.

\proofstep{5. Away-inertia invariants and the original uniformizer.}
Away from the coefficient prime, $c_\ell$ and $\kappa$
are unramified.  In the first part the full finite
character has inertia order dividing six, prime to
$\ell$.  In the second part, $\xi_7^a$ is unramified
away seven, while $\chi^j\beta_d$ has inertia order
dividing six, again prime to seven.
A nontrivial inertia value has unit difference from
one, so $D^{I_v}=0$; if the inertia action is trivial,
$D^{I_v}=D$, which is $\varpi$-divisible.
This includes every place above $2,3,7,269$ that is
away from the coefficient prime, as well as all others.
Hence all local invariant quotients modulo $\varpi$ vanish.

The local and global long exact sequences of
$0\to A\to D\xrightarrow{\varpi}D\to0$ contain
\[
 D^G/\varpi D^G\longrightarrow H^1(G,A)
          \longrightarrow H^1(G,D)[\varpi]\longrightarrow0.
 \tag{MQS.18}
\]
A global strict class killed by $\varpi$ has a residual
lift.  Every local restriction of that lift is zero,
because its image is zero and the local invariant
quotient is zero.  The lift is full strict, hence zero
by Step3.  The original strict group has no
$\varpi$-torsion.  Continuous cocycles into the discrete
$E/O$ have finite image and are killed by a power of
$\varpi$; any nonzero strict class would produce
nonzero $\varpi$-torsion.  This proves (MQS.6) for the
unramified-prime and order-two cases on the original $O$,
without an a priori finiteness assumption.

\proofstep{6. The entire residual-trivial seven-factor family.}
It remains to prove the matched case $\rho_{a,e}=1$,
$\beta_d\ne1$, at coefficient seven.  The fully proved
AC-G211 applies to the three original quadratic factors
$\chi_{-4},\chi_8,\chi_{-8}$, positive weight $e'=6$,
and any continuous near-one character unramified away seven.
Define, on the literal original coefficient ring,
\[
                \kappa'=\xi_7^{-a}\kappa c_7^{e-6}.
 \tag{MQS.19}
\]
Since $\rho_{a,e}=1$ and $\bar c_7^6=1$,
$\bar\kappa'=1$.  All factors are unramified away seven,
so $\kappa'$ has exactly the permitted domain.  The full
original coefficient, not just its reduction, satisfies
\[
 (E/O)(\chi^j\beta_d(\kappa')^{-1}c_7^{-6})
          =(E/O)(\chi^j\beta_d\xi_7^a\kappa^{-1}c_7^{-e}).
 \tag{MQS.20}
\]
AC-G211 therefore proves both strict vanishing statements
for this entire case, with the identical completion and
away-inertia kernels.  No redefinition of the actual
weight, finite label or coefficient order is made in the
application; the auxiliary $e'=6$ is only the parameter
in the already proved all-near-one-character theorem.
The case $\beta_d=1$, $\rho_{a,e}=1$ is not included.

\proofstep{7. The full actual family, counts and unchanged remainder.}
The original host (FG.3) contains $K,t,\zeta_7,\zeta_8$.
The choices of two-power character $\beta_d$, rational
seven-character and Kummer power vary independently
in (FC.1).  All labels in (MQS.8) are
therefore literal distinct original finite labels;
their auxiliary $q,d$ factors are trivial.  The
squarefree integer $d$ here is not the auxiliary
growing source parameter denoted by the same letter.
Their values lie in $\mu_6$, already contained in
each prescribed original coefficient completion.
Set $\kappa=1$.  Inflation from the original global
ramification quotient into $G_K$ is injective and
preserves the prescribed completion and inertia
kernels.  Thus the actual full strict groups inject
into (MQS.6), with the same $O_\ell$, embedding,
uniformizer and rank $b_\ell$.

For fixed $d$ and prime, both Kummer powers are
admitted at exactly the positive multiples of
$\ell-1$.  Their number below integer $T$ proves
the first count in (MQS.9).  At coefficient seven,
$\bar\xi_7^a$, $a=0,\ldots,5$, enumerate all six
rational characters of conductor dividing seven
in each original embedding.  Multiplication by
$\bar c_7^{-e}$ permutes these six characters.
Exactly one product has order two and exactly one
is trivial.  Four two-power factors in the first
case and three in the second, with both Kummer
powers, give $8+6=14$ labels per positive weight.
In the order-two case $a=0$ occurs exactly when
$e\equiv3\pmod6$, for all four factors; in the
matched case it occurs exactly when $6\mid e$.

The unramified-prime rows have a nontrivial rational
quadratic seven-factor, so they are disjoint from
AC-G203, AC-G207 and AC-G215 and from the even
rational seven-powers of AC-G209.
At coefficient seven both admitted residual rational
orders are different from AC-G213's order three.
The AC-G211 rows have $\beta_d\ne1$, $6\mid e$ and
$a=0$; they are precisely the old overlap with
the matched case, giving $6\lfloor N/6\rfloor$ rows.
There is no overlap with the new order-two case.
The two sets also have different coefficient
primes.  Their entire (CP.3) mixed contributions
are $\log1/b_\ell=0$.  Deleting them from any
retained $J$ proves (MQS.10), without changing the
rational noncritical sum, adding a comparison cost,
or choosing another window.
\end{proof}

\paragraph{Finite-cover dictionary, sources and ownership.}
The original $D[\varpi^n]$ are finite locally constant
etale sheaves on
$\operatorname{Spec}\mathcal O_K[1/(2\cdot3\cdot7\cdot269\cdot\ell)]$.
For the unramified-prime and order-two cases, the residual
action is trivialized by the exact $C_3\times C_2$ cover
$F_d/K$.  For the matched case, the corresponding
retained G211 cover is $L(i)$, $L(\sqrt2)$ or $L(\sqrt{-2})$
according to the original $\beta_d$.
Full completion and inertia restriction kernels are retained literally.
This is the existing SGA~1, Expose V, sections 4 and 7
finite-cover/Artin--Tate dictionary at the audited
(FC.1)--(FC.2) scope, with no new archival or motive claim.
The class-field and norm statements retain the exact
complete applied Sharifi scopes; the natural class
generator contract retains the applied PARI
factorbase/Smith and prime-test scopes from G211/G213.
Root proposed the full four-field route and owns
its class and first action computations.  Spectral
derives the entire two-scope original-coefficient
proof and counts.  Trace's independently generated
certificates and rational ideal checks are separate
evidence with their own execution scope.  The unused
degree-eighteen full-rational-seven experiment is not
an input to this theorem.  Neither its pending status
nor later packaging supplies another arithmetic claim.
